\documentclass[twoside,10pt]{../jmlr} 
% \documentclass{colt2020} % Include author names

% The following packages will be automatically loaded:
% amsmath, amssymb, natbib, graphicx, url, algorithm2e

\allowdisplaybreaks
\newcommand{\jmlrBlackBox}{\rule{1.5ex}{1.5ex}}
\providecommand{\BlackBox}{\jmlrBlackBox}
\newcommand{\jmlrQED}{\hfill\jmlrBlackBox\par\bigskip}
\usepackage{graphicx}
\makeatletter
 \let\Ginclude@graphics\@org@Ginclude@graphics 
\makeatother

\newcommand{\rb}[1]{\textcolor{magenta}{RB:#1}}

%\pdfoutput=1

%  \newcommand{\jmlrBlackBox}{\rule{1.5ex}{1.5ex}}
%    \providecommand{\BlackBox}{\jmlrBlackBox}
%  \newcommand{\jmlrQED}{\hfill\jmlrBlackBox\par\bigskip}


%\allowdisplaybreaks

\title{grifting by optimal transport}
\usepackage{times}
\usepackage[utf8]{inputenc} % allow utf-8 input
\usepackage[T1]{fontenc}    % use 8-bit T1 fonts
\usepackage{hyperref}       % hyperlinks
\usepackage{url}            % simple URL typesetting
\usepackage{booktabs}       % professional-quality tables
\usepackage{amsfonts}       % blackboard math symbols
\usepackage{nicefrac}       % compact symbols for 1/2, etc.
\usepackage{microtype}      % microtypography
\usepackage{xspace}
\usepackage{amsmath}


\usepackage{mathtools}  
\usepackage{amsmath}
\usepackage{amssymb}
\usepackage{tabulary}
\usepackage{booktabs}

\usepackage{bm}


\newtheorem{assumption}{Assumption}
\usepackage{algorithm}
\usepackage{algpseudocode}
\usepackage{algpascal}
\usepackage{comment}
\usepackage{selectp}
\usepackage{selectp}





%\newcommand{\dirac}[2]{\delta_{#1}(#2)}
\newcommand{\kprimal}{K}
\newcommand{\dirac}[1]{\delta_{#1}}
\newcommand{\ustar}{u^\star}
\newcommand{\vstar}{v^\star}
\newcommand{\psistar}{\psi^\star}
\newcommand{\Vstar}{V^\star}
\newcommand{\Xstar}{X^\star}
\newcommand{\Ystar}{Y^\star}
\newcommand{\Vhat}{\wh{V}}
\newcommand{\pistar}{\pi^\star}
\newcommand{\pibar}{\overline{\pi}}
\newcommand{\mustar}{\mu^\star}
\newcommand{\hmu}{\wh{\mu}}

\newcommand{\Hplus}{H^+}
\newcommand{\Htest}{H_{\text{test}}}

\newcommand{\muout}{\mu_{\text{out}}}

\newcommand{\mupi}{\mu^\pi}
\newcommand{\nupi}{\nu^\pi}
\newcommand{\dpi}{d^\pi}
\newcommand{\hy}{\wh{y}}
\newcommand{\hx}{\wh{x}}

\newcommand{\piout}{\pi_{\text{out}}}

\newcommand{\tPi}{\wt{\Pi}}
\newcommand{\tpi}{\wt{\pi}}
\newcommand{\tmu}{\wt{\mu}}
\newcommand{\tc}{\wt{c}}
\newcommand{\bmu}{\overline{\mu}}
\newcommand{\bmuX}{\overline{\mu}_\X}
\newcommand{\HH}[2]{\mathcal{H}\pa{#1\middle\|#2}}
\newcommand{\HHX}[2]{\mathcal{H}_{\X}\pa{#1\middle\|#2}}

\newcommand{\bone}{\bm{1}}
\newcommand{\diag}{\mathrm{diag}}
\newcommand{\dd}{\mathrm{d}}
\newcommand{\nux}{\nu_\X}
\newcommand{\nuy}{\nu_\Y}

\newcommand{\bPi}{\overline{\Pi}}

\newcommand{\X}{\mathcal{X}}
\newcommand{\U}{\mathcal{U}}
\newcommand{\XI}{\X^{\infty}}
\newcommand{\PX}{\mathcal{P}_X}
\newcommand{\PY}{\mathcal{P}_Y}
\newcommand{\Y}{\mathcal{Y}}
\newcommand{\YI}{\mathcal{Y}^{\infty}}
\newcommand{\V}{\mathcal{V}}
\newcommand{\LL}{\mathcal{L}}
\newcommand{\GG}{\mathcal{G}}
\newcommand{\hGG}{\wh{\GG}}
\newcommand{\BB}{\mathcal{B}}
\newcommand{\TT}{\mathcal{T}}
\newcommand{\TTpi}{\mathcal{T}^{\pi}}
\newcommand{\TTpik}{\mathcal{T}^{\pi_k}}
\newcommand{\TTP}{\mathcal{T}_{P}}
\newcommand{\TTQ}{\mathcal{T}_{Q}}
\newcommand{\TTX}{\mathcal{T}_{\X}}
\newcommand{\TTY}{\mathcal{T}_{\Y}}
\newcommand{\BX}{\mathcal{B}_\mathcal{X}}
\newcommand{\BY}{\mathcal{B}_\mathcal{Y}}
\newcommand{\SX}{\mathcal{S}_\mathcal{X}}
\newcommand{\SY}{\mathcal{S}_\mathcal{Y}}
\newcommand{\PPX}{P_\mathcal{X}}
\newcommand{\PPY}{P_\mathcal{Y}}
\newcommand{\nuzerox}{\nu_{0,\X}}
\newcommand{\nuzeroy}{\nu_{0,\Y}}


\newcommand{\F}{\mathcal{F}}
\newcommand{\B}{\mathcal{B}}
\newcommand{\FX}{\mathcal{F}_{\mathcal{X}}}
\newcommand{\FXI}{\mathcal{F}_{\mathcal{X}}^{\infty}}
\newcommand{\FY}{\mathcal{F}_{\mathcal{Y}}}
\newcommand{\FYI}{\mathcal{F}_{\mathcal{Y}}^{\infty}}
% \newcommand{\A}{\mathcal{A}}
\newcommand{\N}{\mathcal{N}}
% \newcommand{\I}{\mathcal{I}}
\newcommand{\real}{\mathbb{R}}
\newcommand{\Dw}{\mathcal{D}}
\newcommand{\Sw}{\mathcal{S}}
\newcommand{\Pw}{\mathcal{P}}
\newcommand{\Rw}{\mathcal{R}}
\newcommand{\Ww}{\mathcal{W}}
\newcommand{\DD}[2]{\mathcal{D}\pa{#1\middle\|#2}}
\newcommand{\DDh}[2]{\mathcal{B}_{\bH}\pa{#1\middle\|#2}}
\newcommand{\DDp}[2]{\mathcal{D}_{\ell_p}\pa{#1\middle\|#2}}
\newcommand{\DDLC}[2]{\mathcal{D}_{\mathrm{LC}}\pa{#1\middle\|#2}}
\newcommand{\tDD}[2]{\wt{\mathcal{D}}\pa{#1\middle\|#2}}
\newcommand{\DDR}[2]{\mathcal{D}_R\pa{#1\middle\|#2}}
\newcommand{\DDsigma}[2]{\mathcal{D}_{\sigma}\pa{#1\middle\|#2}}
\newcommand{\DDKL}[2]{\mathcal{D}_{\mathrm{KL}}\pa{#1\middle\|#2}}
\newcommand{\DDxy}[2]{\mathcal{D}_{\mathrm{KL}}^{xy}\pa{#1\middle\|#2}}
\newcommand{\kldiv}[2]{d_{\mathrm{kl}}\pa{#1\middle\|#2}}
\newcommand{\DDPhi}[2]{\mathcal{B}_\Phi\pa{#1\middle\|#2}}
\newcommand{\DDalpha}[2]{\mathcal{D}_{\alpha}\pa{#1\middle\|#2}}
\newcommand{\DDpsi}[2]{\mathcal{D}_{\psi}\pa{#1\middle\|#2}}
\newcommand{\DDchi}[2]{\mathcal{D}_{\chi^2}\pa{#1\middle\|#2}}
\newcommand{\DDphi}[2]{\mathcal{D}_{\varphi}\pa{#1\middle\|#2}}
\newcommand{\DDPhiS}[2]{\mathcal{D}_{\Phi_S}\pa{#1\middle\|#2}}
\newcommand{\DDtPhi}[2]{\mathcal{D}_{\tPhi}\pa{#1\middle\|#2}}
\newcommand{\DDS}[2]{\mathcal{D}^S\pa{#1\middle\|#2}}
\newcommand{\DDc}[3]{\mathcal{D}^{#3}\pa{#1\middle\|#2}}
\newcommand{\tDDS}[2]{\wt{\mathcal{D}}^S\pa{#1\middle\|#2}}
\newcommand{\DDf}[2]{\mathcal{D}_\varphi\pa{#1\middle\|#2}}
\newcommand{\OO}{\mathcal{O}}
\newcommand{\tOO}{\wt{\OO}}
\newcommand{\trace}[1]{\mbox{trace}\left(#1\right)}
\newcommand{\II}[1]{\mathbb{I}_{\left\{#1\right\}}}
\newcommand{\PP}[1]{\mathbb{P}\left[#1\right]}
\newcommand{\PPZ}[1]{\mathbb{P}_Z\left[#1\right]}
\newcommand{\EE}[1]{\mathbb{E}\left[#1\right]}
\newcommand{\EEP}[1]{\mathbb{E}_P\left[#1\right]}
\newcommand{\EES}[1]{\mathbb{E}_S\left[#1\right]}
\newcommand{\EESb}[1]{\mathbb{E}_S\bigl[#1\bigr]}
\newcommand{\EESB}[1]{\mathbb{E}_S\Bigl[#1\Bigr]}
\newcommand{\EEZ}[1]{\mathbb{E}_Z\left[#1\right]}
\newcommand{\EEb}[1]{\mathbb{E}\bigl[#1\bigr]}
\newcommand{\EEtb}[1]{\mathbb{E}_t\bigl[#1\bigr]}
\newcommand{\EXP}{\mathbb{E}}
\newcommand{\EEpi}[1]{\mathbb{E}_{\pi}\left[#1\right]}
\newcommand{\PPpi}[1]{\mathbb{P}_{\pi}\left[#1\right]}
\newcommand{\EEs}[2]{\mathbb{E}_{#2}\left[#1\right]}
\newcommand{\EEu}[1]{\mathbb{E}_{\u}\left[#1\right]}
\newcommand{\EEst}[1]{\mathbb{E}_{*}\left[#1\right]}
\newcommand{\EEo}[1]{\mathbb{E}_{0}\left[#1\right]}
\newcommand{\PPt}[1]{\mathbb{P}_t\left[#1\right]}
\newcommand{\EEt}[1]{\mathbb{E}_t\left[#1\right]}
\newcommand{\PPi}[1]{\mathbb{P}_i\left[#1\right]}
\newcommand{\EEi}[1]{\mathbb{E}_i\left[#1\right]}
\newcommand{\PPc}[2]{\mathbb{P}\left[#1\left|#2\right.\right]}
\newcommand{\PPs}[2]{\mathbb{P}_{#2}\left[#1\right]}
\newcommand{\PPcs}[3]{\mathbb{P}_{#3}\left[#1\left|#2\right.\right]}
\newcommand{\PPct}[2]{\mathbb{P}_t\left[#1\left|#2\right.\right]}
\newcommand{\PPcc}[2]{\mathbb{P}\left[\left.#1\right|#2\right]}
\newcommand{\PPcct}[2]{\mathbb{P}_t\left[\left.#1\right|#2\right]}
\newcommand{\PPcci}[2]{\mathbb{P}_i\left[\left.#1\right|#2\right]}
\newcommand{\EEc}[2]{\mathbb{E}\left[#1\left|#2\right.\right]}
\newcommand{\EEcc}[2]{\mathbb{E}\left[\left.#1\right|#2\right]}
\newcommand{\EEcs}[3]{\mathbb{E}_{#3}\left[\left.#1\right|#2\right]}
\newcommand{\EEcct}[2]{\mathbb{E}_t\left[\left.#1\right|#2\right]}
\newcommand{\EEcci}[2]{\mathbb{E}_i\left[\left.#1\right|#2\right]}
%\renewcommand{\th}{\ensuremath{^{\mathrm{th}}}}
\def\argmin{\mathop{\mbox{ arg\,min}}}
\def\argmax{\mathop{\mbox{ arg\,max}}}
\newcommand{\ra}{\rightarrow}

\newcommand{\siprod}[2]{\langle#1,#2\rangle}
\newcommand{\iprod}[2]{\left\langle#1,#2\right\rangle}
\newcommand{\Siprod}[2]{\left\langle#1,#2\right\rangle_{|S}}
\newcommand{\biprod}[2]{\bigl\langle#1,#2\bigr\rangle}
\newcommand{\Biprod}[2]{\Bigl\langle#1,#2\Bigr\rangle}
\newcommand{\norm}[1]{\left\|#1\right\|}
\newcommand{\bnorm}[1]{\bigl\|#1\bigr\|}
\newcommand{\onenorm}[1]{\norm{#1}_1}
\newcommand{\twonorm}[1]{\norm{#1}_2}
\newcommand{\infnorm}[1]{\norm{#1}_\infty}
\newcommand{\tvnorm}[1]{\norm{#1}_{\mathrm{TV}}}
\newcommand{\Hnorm}[1]{\norm{#1}_{\mathcal{H}}}
\newcommand{\ev}[1]{\left\{#1\right\}}
% \newcommand{\abs}[1]{\left|#1\right|}
\newcommand{\babs}[1]{\bigl|#1\bigr|}
\newcommand{\pa}[1]{\left(#1\right)}
\newcommand{\bpa}[1]{\bigl(#1\bigr)}
\newcommand{\Bpa}[1]{\Bigl(#1\Bigr)}
\newcommand{\BPA}[1]{\Biggl(#1\Biggr)}
\newcommand{\sign}{\mbox{sign}}

\newcommand{\wh}{\widehat}
\newcommand{\wt}{\widetilde}


\newcommand{\hr}{\wh{r}}
\newcommand{\hf}{\wh{f}}
\newcommand{\hs}{\wh{s}}
\newcommand{\bA}{\overline{A}}
\newcommand{\deff}{d_{\text{eff}}}

\newcommand{\Cinf}{C_\infty}
\newcommand{\tfs}{\wt{f}_{\sigma}}
\newcommand{\tQ}{\wt{Q}}
\newcommand{\tW}{\wt{W}}
\newcommand{\tR}{\wt{R}}
\newcommand{\tPhi}{\wt{\Phi}}
\newcommand{\tZ}{\wt{Z}}
\newcommand{\tS}{\wt{S}}
\newcommand{\tD}{\wt{D}}
\newcommand{\tw}{\wt{w}}
\newcommand{\bw}{\overline{w}}
\newcommand{\hX}{\wh{X}}
\newcommand{\hS}{\wh{\Sigma}}
\newcommand{\hw}{\wh{w}}
\newcommand{\bX}{\overline{X}}
\newcommand{\bY}{\overline{Y}}
\newcommand{\bS}{\overline{S}}
\newcommand{\bZ}{\overline{Z}}
\newcommand{\bW}{\overline{W}}
\newcommand{\bP}{\overline{P}}
\newcommand{\bDelta}{\overline{\Delta}}
\newcommand{\Pj}{P_X\otimes P_Y}
\newcommand{\barf}{\overline{f}}
\newcommand{\hQ}{\wh{Q}}
\newcommand{\hP}{\wh{P}}
\newcommand{\tP}{\wt{P}}
\newcommand{\transpose}{^\mathsf{\scriptscriptstyle T}}

\newcommand{\bloss}{\overline{\ell}}
\newcommand{\hloss}{\wh{\ell}}
\newcommand{\loss}{\ell}
\newcommand{\gen}{\textup{gen}}

\usepackage{todonotes}
\definecolor{PalePurp}{rgb}{0.66,0.57,0.66}
\newcommand{\todoG}[1]{\todo[color=PalePurp!30]{\textbf{Grg:} #1}}
\newcommand{\todoGI}[1]{\todo[inline,color=PalePurp!30]{#1}}
% % \newcommand{\reddd}[1]{}
\newcommand{\redd}[1]{\textcolor{red}{#1}}
\newcommand{\grg}[1]{\textcolor{red}{[\textbf{Grg:} #1]}}
\newcommand{\lud}[1]{\textcolor{red}{[\textbf{Ludovic:} #1]}}
\newcommand{\ser}[1]{\textcolor{red}{[\textbf{Sergio:} #1]}}
\newcommand{\gbr}[1]{\textcolor{orange}{[\textbf{Gbr:} #1]}}

\newcommand{\regret}{\mathrm{regret}}

% \newcommand{\qed}{\hfill\BlackBox\\[2mm]}


\newcommand{\hL}{\wh{L}}
\newcommand{\snr}{\mathtt{SNR}}

\newcommand{\alg}{\mathcal{A}}
\newcommand{\Zw}{\mathcal{Z}}
\newcommand{\bL}{\overline{L}}
%\newcommand{\bH}{\overline{H}}
\newcommand{\bH}{h}
\renewcommand{\th}{\wt{h}}
\newcommand{\Law}{\text{Law}}

\usepackage[normalem]{ulem}
\usepackage{enumitem}

% Put your macros here

\newcommand\am[1]{{\color{red}(\textbf{AM:} #1)}}

\newcommand\pSrc{p_{\text{src}}}
\newcommand\Pm{\mathbb{P}}
\newcommand\supp{\text{supp}}
\newcommand\opt{\text{OPT}}

\newcommand\src{\nu_{\text{src}}}
\newcommand\tgt{\nu_{\text{tgt}}}
\newcommand\dSrc{\rho_{\text{src}}}
\newcommand\dTgt{\rho_{\text{tgt}}}
\newcommand\Wc{\mathcal{W}}
\newcommand\fwd{\mathcal{F}}
\newcommand\bwd{\mathcal{B}}
\newcommand\Sc{\mathcal{S}}
% \newcommand\norm[1]{\left\|#1\right\|}
\newcommand\Lc{\mathcal{L}}
\newcommand\hLc{\wh{\Lc}}
%\newcommand\Cc{\mathcal{C}}

% \usepackage{etoc}
% \usepackage{authblk}
% \usepackage{microtype}
% \usepackage{graphicx}
% %\usepackage{subcaption}
% \usepackage[T1]{fontenc}
% \usepackage[english]{babel}
% \usepackage[utf8]{inputenc}
% \usepackage{booktabs}
% \usepackage{array}
% \usepackage{amsmath}
% \usepackage{amssymb}
% \usepackage{mathtools}
% \usepackage{amsthm}
% \usepackage{bm}
% \usepackage{thmtools}
% \usepackage{thm-restate}
% \usepackage{hyperref}
% \hypersetup{
%     colorlinks = true,
%     linkcolor = blue,
%     citecolor=blue,
%     urlcolor=black,
% }
% \usepackage{xcolor}
% \definecolor{mblue}{rgb}{0.0,0.0,1.0}
% \definecolor{mgreen}{rgb}{0.0,0.501960784314,0.0}
% \definecolor{mred}{rgb}{1.0,0.0,0.0}
% \usepackage[capitalize]{cleveref}
% \usepackage[ruled]{algorithm2e}
% \usepackage[noend]{algpseudocode}
% \renewcommand{\algorithmiccomment}[1]{\hfill$\triangleright~${\color{gray}#1}}
% \usepackage{bbm}
% \usepackage{xcolor}
% \usepackage{parskip}
% \usepackage{setspace}
% %\usepackage{etoolbox}
% \usepackage{comment}
% \usepackage{caption}
% \usepackage{enumitem}
% \usepackage{wrapfig}
% 
% \newtheorem{theorem}{Theorem}[section]
% \newtheorem{example}{Example}[section]
% \newtheorem{remark}{Remark}[section]
% \newtheorem{corollary}{Corollary}[section]
% \newtheorem{definition}{Definition}[section]
% \newtheorem{lemma}{Lemma}[section]
% \newtheorem{claim}{Claim}[section]
% \newtheorem{observation}{Observation}[section]
% \newtheorem{proposition}{Proposition}[section]
% \newtheorem{assumption}{Assumption}[section]
% \newtheorem{question}{Question}
% \newtheorem{itheorem}{Informal Theorem}[section]

\newcommand\R{\mathbb{R}}
\newcommand\inner[2]{\left\langle #1, #2 \right\rangle}
\newcommand\indicator[1]{\mathbbm{1}\left\{ #1 \right\}}
\newcommand\Simplex[1]{\Delta_{#1}}

\newcommand\Ac{\mathcal{A}}
\newcommand\Bc{\mathcal{B}}
\newcommand\Cc{\mathcal{C}}
\newcommand\Dc{\mathcal{D}}
\newcommand\Fc{\mathcal{F}}
\newcommand\Xc{\mathcal{X}}
\newcommand\Yc{\mathcal{Y}}
\newcommand\Mc{\mathcal{M}}
\newcommand\Tc{\mathcal{T}}
%\newcommand\E{\mathbb{E}}
\DeclareMathOperator*{\E}{\mathbb{E}}
\newcommand\Dkl{D_{\text{KL}}}
\newcommand\pre{p^{\text{pre}}}
\newcommand\dHat{\hat{d}}
\newcommand\sHat{\hat{s}}
\newcommand\pHat{\hat{p}}

\newcommand\OTilde{\tilde{O}}
\newcommand\biggVert{~\bigg\vert~}
\newcommand\bigVert{~\big\vert~}
\newcommand\normVert{~\vert~}
\newcommand\rHat{\hat{r}}
\newcommand\floor[1]{\lfloor #1 \rfloor}
\newcommand\ceil[1]{\lceil #1 \rceil}
\newcommand{\bc}[1]{\left\{{#1}\right\}}
\newcommand{\br}[1]{\left({#1}\right)}
\newcommand{\bs}[1]{\left[{#1}\right]}
\newcommand{\abs}[1]{\left\lvert #1 \right\rvert}
\newcounter{protocol}
\makeatletter
\newenvironment{protocol}[1][htb]{%
  \let\c@algorithm\c@protocol
  \renewcommand{\ALG@name}{Protocol}
  \begin{algorithm}[#1]%
  }{\end{algorithm}
}

\makeatletter
\providecommand*{\bigcupdot}{%
  \mathop{%
    \vphantom{\bigcup}%
    \mathpalette\@bigcupdot{}%
  }%
}
\newcommand*{\@bigcupdot}[2]{%
  \ooalign{%
    $\m@th#1\bigcup$\cr
    \sbox0{$#1\bigcup$}%
    \dimen@=\ht0 %
    \advance\dimen@ by -\dp0 %
    \sbox0{\scalebox{1.5}{$\m@th#1\cdot$}}%
    \advance\dimen@ by -\ht0 %
    \dimen@=.5\dimen@
    \hidewidth\raise\dimen@\box0\hidewidth
  }%
}
\makeatother



\author{
 \Name[{Gergely~Neu}]{Gergely Neu} \Email{gergely.neu@gmail.com}\\
 \addr Universitat Pompeu Fabra, Barcelona, Spain
 \AND
 \Name{(add-your-name)} \Email{add-your-email@gmail.com}\\
 \addr University of XYZ
 }
 
\begin{document}

\maketitle

\begin{abstract}
aaaauuuuu
\end{abstract}

\begin{keywords}%
  optimal transport, generative modeling, Markov chains
\end{keywords}


\section{Problem formulation (2026.05.15.)}
``Drifting models'' is a new (and somehow very popular) approach for generative modeling for learning 
to sample from a given target distribution $\tgt$ accessible only via samples, given samples from a source distribution 
$\src$ as input. Essentially, the idea behind drifting models is to define a Markov chain whose stationary distribution 
is equal to the target distribution, and directly estimate the stationary distribution of the proposed chain. The chain 
considered in the original work is a simple process that operates via rules of attraction and repellence: roughly 
speaking, a state is repelled by the source distribution and attracted by the target distribution. The original paper 
argues that this process converges to the target distribution under some conditions (once the specifics of the policy 
are set up correctly). Is there a good reason to believe that the proposed process is the best possible one for 
purposes of generative modeling? This is the question we ask in this note.

We thus formulate our problem as finding a control policy over the state space $\X = \real^d$ that yields $\tgt$ as its 
stationary distribution. There are many possible policies that satisfy this, e.g., the process above or the one that 
draws states i.i.d.~from the target distribution, and so on. We try to formulate criteria for a good policy in the 
language of infinite-horizon undiscounted optimal control below. The main idea is to find a policy with minimal 
\emph{transient cost} incurred before the stationary distribution is reached. For concreteness, let $\pi:\X \ra 
\Delta_{\X}$ be a stochastic policy and let us consider its stationary distribution $\nupi$ (assuming that it exists). 
We denote measures on $\X$ by $\nu$ and the action of the transition kernel $\pi$ is taking $\nu$ to the distribution 
$\pi\nu$ with density $(\pi\nu)(x) = \int_{\X} \pi(x|x^-) \nu(x^-) \dd x^-$ (assuming everything has densities). Then, 
the transient state distribution induced by $\pi$ (under initial-state distribution $\src$) is defined as
\begin{align*}
 \dpi &= \sum_{t=0}^\infty \pa{\pi^t \src - \nupi} = \sum_{t=0}^\infty \pi^t \pa{\src - \nupi}
 \\
 &= \src - \nupi + \sum_{t=1}^\infty \pi^t \pa{\src - \nupi} = \src - \nupi + \pi \dpi.
\end{align*}
Well, the last part is not part of the definition but rather a consequence of the definition, which can be seen as the 
equivalent of the classic flow constraints for occupancy measures. We thus refer to $\dpi$ as the \emph{transient 
occupancy measure} of $\pi$.

After reparametrizing the problem as in the classic LP formulations to eliminate the product $\pi\dpi$, we thus define 
the state-next-state transient occupancy $\mupi(x,x')=\dpi(x)\pi(x'|x)$ instead and note that it satisfies the linear 
constraints
\[
 B\mupi = \src - \nupi + F\mupi,
\]
where we have defined the forward and backward-marginalization operators $F$ and $B$ with their respective actions 
$(F\mu)(x) = \int_{\X} \mu(x',x) \dd x'$ and $(B\mu)(x) = \int_{\X} \mu(x,x') \dd x'$. Similarly, the 
stationary distribution $\nupi$ satisfies
\[
 \nupi = \pi \nupi.
\]
This one still has a $\pi$ in there which is no good. How can we reconcile this with the state-next-state occupancies? 
Those should obey the same conditional distribution $\pi$, which seems difficult to extract from $\mu$ while retaining 
the linearity of the constraints.

Here's an idea that might work. Since we hope to find a policy with stationary distribution equal to $\tgt$, we might 
as well replace the desired stationary distribution with this measure (and hope that the solution will somehow indeed 
yield a policy with this limiting measure). We thus simply consider the following linear program:
\begin{align*}
 \min_{\mu} \int_{\X} \int_{\X} &\mu(x,x') \norm{x-x'}^2 \dd x \dd x'
 \\
 \mbox{subject to} \qquad & B\mu = \src - \tgt + F\mu.
\end{align*}
LOL this just looks like the classic Kantorovich LP all over again, except with the two constraints folded into one! 
Can this still be a good idea somehow? Let us consider a feasible point $\mu$ and try to understand it as a transient 
occupancy measure (denoting $\mu = \mupi$ and $\pi$ as the corresponding kernel):
\[
 \dpi = \src - \tgt + \pi\dpi.
\]
Indeed, on the left-hand side we have $\dpi(x) = \int_{\X} \mupi(x,x') \dd x'$ by definition, and on the right-hand 
side we have $\int_{\X} \mupi(x,x') \dd x = \int_{\X} \dpi(x) \pi(x'|x) \dd x$. Thus, we can iterate the above relation 
and get
\[
 \dpi = \src - \tgt + \pi\dpi = \src - \tgt + \pi(\src - \tgt) + \pi^2 \dpi = \dots = \sum_{t=0}^\infty \pi^t\pa{\src - 
\tgt}.
\]
This sum is convergent if $\pi$ has appropriate mixing properties. But notice that it does not imply that $\tgt$ would 
be its stationary distribution. So we have basically proved that any quickly-mixing policy $\pi$ is feasible, and that 
$\dpi$ corresponds to the \emph{difference} of transient measures when started from distributions $\src$ and $\tgt$.

So the problem remains: we should find a formulation that does actually yield $\tgt$ as the stationary distribution. 
Maybe just explicitly require this via the constraint $F\mu = \tgt$? \grg{NO this is silly: $\mu$ is not directly 
related to the stationary distribution with equations like that.} This of course brings us back \emph{exactly} to the 
Kantorovich LP (once the first constraint is simplified appropriately), which also verifies that the optimal policy 
should be such that it maps $x$ deterministically to $M^\star(x)$ using the Monge map $M^\star$. But this doesn't mean 
that the above reformulation as a sequential decision-making problem is useless! We could in principle still use ideas 
from the VDT paper to solve this control problem via primal-dual methods: characterize the value functions by a variant 
of the Bellman equations (Poisson equations?), the optimal policy in terms of the gradient of the value function, and 
solve for the optimal value function by primal-dual methods. Maybe it could work, or maybe that's just nonsense... 
should be fun to think about it anyway.

\subsection{Lazy approach 1: cost optimality under stationary distribution}
The most straightforward way to fix up this formulation is to do away with the transient occupancies and simply work 
with stationary distributions. That is, simply consider the following adjustment to the classic undiscounted control LP:
\begin{align*}
 \min_{\mu} \int_{\X} \int_{\X} &\mu(x,x') \norm{x-x'}^2 \dd x \dd x'
 \\
 \mbox{subject to} \qquad & B\mu = F\mu
 \\
 & B\mu = \tgt.
\end{align*}
What is the justification for this objective though? The optimal policy might not be unique here, and there's no way to 
differentiate between ``good'' or ``bad'' ones. The objective only says ``don't move around the states too much once 
you've hit the target'', but doesn't say anything beyond that. (Also, this formulation doesn't say anything about the 
source distribution, which may or may not be a bad thing. If we want to connect with OT, it's a bad thing for sure.)

\subsection{Lazy approach 2: discounted transients}
The second most straightforward approach is to adjust the LP by introducing discounting:
\begin{align*}
 \min_{\mu} \int_{\X} \int_{\X} &\mu(x,x') \norm{x-x'}^2 \dd x \dd x'
 \\
 \mbox{subject to} \qquad & B\mu = (1-\gamma)\src + \gamma F\mu
 \\
 & B\mu = (1-\gamma)\src + \gamma \tgt.
\end{align*}
The source is included in the second constraint to make sure that the problem is feasible. Unlike the previous problem 
formulation, this one is somewhat interpretable: I give you a geometric time horizon to move the source to the target 
with minimal total transport cost. So basically the dynamic OT problem but with a random time horizon. That helps 
because the optimal policy is stationary and the value function does not depend on the time index either; all the 
objects to learn are simpler. It's not clear if the cool properties of the classic dynamic OT solution would be 
inherited as nicely as in the fixed-horizon case... but that might just make things more interesting. Also, note that 
sending $\gamma$ to 1 in the limit recovers the previous problem formulation.










\bibliographystyle{plainnat}
\bibliography{../ngbib}
\end{document}
